% vim: ai sw=2 spell spelllang=cs

\begin{frame}
  \frametitle{Síťové rozhraní}

  \heading{LDD3 kap. 17}

  \bigskip

  \begin{itemize}
  \item Ovladač pro ethX
  \end{itemize}

\end{frame}

\begin{frame}[fragile]{Teoretická síťová vrstva}

\begin{center}
  \tikzstyle{sized} = [minimum height=7mm,minimum width=2.5cm]
  \tikzstyle{tl_statenet} = [tl_state,sized,minimum width=4cm]
  \begin{tikzpicture}[node distance=3mm, font=\footnotesize]
      \node [tl_statenet] (USERA) {User / Application};
      \node [tl_statenet] (APP) [below=of USERA] {Application layer};
      \node [tl_statenet] (TRANS) [below=of APP] {Transport layer};
      \node [tl_statenet] (NET) [below=of TRANS] {Network layer};
      \node [tl_statenet] (LINK) [below=of NET] {Link layer};
      \node [tl_statenet] (HW) [below=of LINK] {Hardware layer};

      \node [sized] (USERAEX) [right=of USERA] {Firefox};
      \node [sized] (APPEX) [right=of APP] {HTTP};
      \node [sized] (TRANSEX) [right=of TRANS] {TCP};
      \node [sized] (NETEX) [right=of NET] {IP};
      \node [sized] (LINKEX) [right=of LINK] {Ethernet driver};
      \node [sized] (HWEX) [right=of HW] {Ethernet};

      \path (HWEX.south east) -- (USERAEX.north east)
	node [midway,rotate=270,anchor=south] {\textbf{Příklady}};
      \path (HW.south west) -- (USERA.north west)
	node [midway,left=4pt,rotate=90,anchor=south] {\textbf{Vrstvy}};

  \end{tikzpicture}
\end{center}

\end{frame}

\begin{frame}[fragile]{Linuxová síťová vrstva}

\begin{center}
  \tikzstyle{sized} = [minimum height=7mm,minimum width=2.5cm]
  \tikzstyle{tl_statenet} = [tl_state,sized,minimum width=4cm]
  \pgfdeclarelayer{background}
  \pgfsetlayers{background,main}
  \begin{tikzpicture}[node distance=3mm, font=\footnotesize]
      \node [tl_statenet] (APP) {Application layer};
      \node [tl_statenet] (SYSCALL) [below=5mm of APP] {System call interface};
      \node [tl_statenet] (PAI) [below=of SYSCALL] {Protocol agnostic interface};
      \node [tl_statenet] (PROTO) [below=of PAI] {Network protocols};
      \node [tl_statenet] (DAI) [below=of PROTO] {Device agnostic interface};
      \node [tl_statenet] (DD) [below=of DAI] {Device drivers};
      \node [tl_statenet] (HW) [below=of DD] {Physical device hardware};

      \node [sized] (USERSPACE) [left=of APP] {\textbf{Userspace}};
      \node [sized] (KERNELSPACE) [left=of PROTO] {\textbf{Kernelspace}};

      \begin{pgfonlayer}{background}
	\node [tl_statenet,fit=(APP)] (US) {};
	\node [tl_statenet,fit=(SYSCALL)(PAI)(PROTO)(DAI)(DD)] (KS) {};
      \end{pgfonlayer}

      \path (US.south east) -- (KS.north east)
	node [midway,inner sep=0pt,minimum size=0pt] (MID1) {};
      \node [left=7cm of MID1] {}
        edge [draw,dashed] (MID1);

      \path (KS.south east) -- (HW.north east)
	node [midway,inner sep=0pt,minimum size=0pt,xshift=2pt] (MID2) {};
      \node [left=7cm of MID2] {}
        edge [draw,dashed] (MID2);

  \end{tikzpicture}

\end{center}

\end{frame}

\begin{frame}
  \frametitle{Síťové rozhraní ovladače}

  \heading{Ovladač síťové karty}
  \begin{itemize}
    \item Mezivrstva mezi HW a Linuxem (síťovou vrstvou)
    \item Vytváří \code{eth*} a podobná zařízení (tzv.\ \code{netdevice})
    \item Obsahuje háčky
    \begin{itemize}
      \item Zařízení je UP, DOWN, změna MTU, pošli paket, \dots
    \end{itemize}
  \end{itemize}

  \pause
  \medskip

  \heading{API}
  \begin{itemize}
    \item \header{linux/netdevice.h}, \code{struct net\_device}
    \item Alokace: \code{alloc_netdev}$_D$, \code{free_netdev}$_D$
    \item Registrace: \code{register_netdev}, \code{unregister_netdev}
    \item Háčky: \code{struct net_device_ops} (podobně jako znaková zařízení)
  \end{itemize}

\end{frame}

\begin{frame}[fragile]
  \frametitle{\code{struct net_device_ops}}

  \begin{block}{\code{struct net_device_ops}}
  \begin{lstlisting}
  int (*ndo_open)(struct net_device *dev); /* ip link set up */
  int (*ndo_stop)(struct net_device *dev); /* ip link set down */

  /* send a packet */
  netdev_tx_t (*ndo_start_xmit)(struct sk_buff *skb, struct net_device *dev);

  /* change MTU */
  int (*ndo_change_mtu)(struct net_device *dev, int new_mtu);

  ... /* no receive packet */
  \end{lstlisting}
  \end{block}

\end{frame}

\begin{frame}
  \frametitle{Konkrétnější rozhraní}

  \begin{itemize}
    \item Místo \code{alloc_netdev} se použije konkrétnější \code{alloc_*}
    \begin{itemize}
      \item Inicializuje \code{netdevice}
      \item Délku hlavičky, adresy, fronty, \dots
    \end{itemize}

    \pause

    \item Příklady
    \begin{itemize}
      \item \emph{Ethernet}: \code{alloc_etherdev}
	(\header{linux/etherdevice.h})
      \item IRDA: \code{alloc_irdadev} (\header{net/irda/irda\_device.h})
      \item CAN: \code{alloc_candev} (\header{linux/can/dev.h})
      \item HDLC: \code{alloc_hdlcdev} (\header{linux/hdlc.h})
      \item \dots
    \end{itemize}
  \end{itemize}

\end{frame}

\begin{frame}
  \frametitle{Vytvoření \code{eth*}}

  \heading{Postup vytvoření ethernetového zařízení v systému}
  \begin{enumerate}
    \item Alokace (\code{alloc_etherdev(priv_size)})
    \begin{itemize}
      \item Alokuje navíc paměť o velikosti \code{priv_size}
      \item \code{priv_size} může být \code{0}
    \end{itemize}

    \pause

    \item Inicializace informací
    \begin{itemize}
      \item Nastavení háčků (\code{net_device_ops} do \code{dev->netdev_ops})

      \pause

      \item Nastavení MAC adresy (\code{dev->dev_addr})

      \pause

      \item Získání ukazatele na paměť alokovanou navíc
      \begin{itemize}
	\item Funkcí \code{netdev_priv(dev)}
	\item Jen při alokaci s \code{priv_size} > 0
      \end{itemize}
    \end{itemize}

    \pause

    \item Registrace (\code{register_netdev})
    \item \dots

    \pause

    \item Deregistrace (\code{unregister_netdev})
    \item Uvolnění (\code{free_netdev})
  \end{enumerate}

\end{frame}

\begin{frame}
  \frametitle{Úkol}

  \heading{Vytvoření \code{eth*}}
  \begin{enumerate}
    \item Dle postupu z předchozího slajdu
    \begin{itemize}
      \item Alokace+registrace v \code{module_init}
      \item Deregistrace+uvolnění v \code{module_exit}
      \item MAC: \code{random_ether_addr}
      \item \code{priv_size}: \code{sizeof(struct timer_list)}
      \begin{itemize}
	\item Nezapomeňte na \code{setup_timer}
      \end{itemize}
    \end{itemize}

    \item Vytvořte \code{open}, \code{stop}, \code{start_xmit}
    \begin{itemize}
      \item Těla vypíšou jen název funkce
      \item \code{start_xmit} vrátí \code{NETDEV_TX_OK}
    \end{itemize}

    \item \emph{NEVKLÁDEJTE} do systému
    \begin{itemize}
      \item Nelze modul odebrat
    \end{itemize}
  \end{enumerate}

\end{frame}

\begin{frame}
  \frametitle{Pakety}

  \heading{Tzv.\ socket buffery}
  \begin{itemize}
    \item \header{linux/skbuff.h}, \code{struct sk_buff}
    \item Obecně: \code{dev_alloc_skb}

    \pause

    \item \emph{Pro netdevice:} \code{netdev_alloc_skb}

    \pause

    \item Uvolnění: \code{dev_kfree_skb}

    \pause

    \item Obsah: \code{skb->data}
    \begin{itemize}
      \item Ukazuje na alokovanou paměť
      \item Plní/čte se klasicky přes \code{memcpy} apod.
    \end{itemize}

    \pause

    \item Délka obsahu: \code{skb->len}

    \pause

    \item Protokol: \code{skb->protocol}
    \begin{itemize}
      \item Při příjmu nutné nastavit před předáním systému!
    \end{itemize}
  \end{itemize}

\end{frame}

\begin{frame}{Pakety}{Odesílání do HW}

  \heading{\code{skb} je parametr \code{start_xmit}}
  \begin{enumerate}
    \item Předat do HW data
    \begin{itemize}
      \item \code{skb->data} o délce \code{skb->len}
    \end{itemize}

    \item Uvolnit \code{skb}
    \begin{itemize}
      \item \code{dev_kfree_skb(skb)}
    \end{itemize}

    \item Vrátit stav
    \begin{itemize}
      \item \code{NETDEV_TX_OK}
      \item \code{NETDEV_TX_BUSY}
    \end{itemize}

    \item Lze pozastavit odesílání
    \begin{itemize}
      \item Při zaplněném HW
      \item \code{netif_stop_queue}
      \item Při volném HW později: \code{netif_wake_queue}
    \end{itemize}
  \end{enumerate}

\end{frame}

\begin{frame}
  \frametitle{Úkol}

  \heading{,,Odesílání'' paketů}
  \begin{enumerate}
    \item Doplňte kód pro odesílání (\code{start_xmit})
    \begin{itemize}
      \item Vypište obsah paketu (\code{print_hex_dump})
      \item Uvolněte \code{skb} (\code{dev_kfree_skb})
    \end{itemize}

    \item Vložte modul do systému
    \item Pusťte si \cmd{arping}
    \begin{itemize}
      \item \cmd{\-I vase_rozhrani nejaka_IP_adresa}
    \end{itemize}

    \item Pozorujte \cmd{dmesg}
  \end{enumerate}

\end{frame}

\begin{frame}{Pakety}{Příjem z HW}

  \heading{Např.\ z přerušení}
  \begin{enumerate}
    \item \code{skb = netdev_alloc_skb(dev, len)}
    \begin{itemize}
      \item \code{len}: např.\ MTU nebo přesná velikost
    \end{itemize}

    \pause

    \item Naplnění daty z HW do \code{skb->data} o délce max.\ \code{len}
    \item Nastavení \code{skb->len}
    \begin{itemize}
      \item \code{skb->len} $\leq$ \code{len}
    \end{itemize}

    \pause

    \item Nastavení protokolu
    \begin{itemize}
      \item \code{skb->protocol = eth_type_trans(skb, dev)}
      \item Je to nutné!
    \end{itemize}

    \pause

    \item Předání \code{skb} jádru
    \begin{itemize}
      \item \code{netif_rx(skb)}
      \item Jádro \code{skb} uvolní
    \end{itemize}

  \end{enumerate}

\end{frame}

\begin{frame}
  \frametitle{Úkol}

  \heading{,,Příjem'' paketů}
  \begin{enumerate}
    \item Emulujte příjem paketů
    \begin{itemize}
      \item V \code{open} nastavte periodický časovač
      \item Časovač zrušte ve \code{stop}
      \item V časovači volejte \code{netif_rx} s paketem z \file{pb173/12/}
      \item Inkrementujte o \code{1} v paketu 24.\ znak
    \end{itemize}

    \item Vložte modul do systému
    \item Spusťte \cmd{tcpdump -ni} na \code{eth} rozhraní
  \end{enumerate}

\end{frame}
